\]
It follows by scaling the unit circle and integrating the power series
for $\log(1-z)$ inside its disk; the boundary case follows from the
integrability of $\log|2\sin(\theta/2)|$. This is also the standard
circle-potential formula \cite{Saff}. In particular
$I(\omega_i)=\log\varepsilon$.

A unit arcsine measure on an interval of length $L$ assigns at most
$4\sqrt{2r}/(\pi\sqrt L)$ to an interval of radius $r$. To see a
sufficient bound, the maximum mass of an interval of fixed length occurs
at an endpoint of the supporting interval; its mass is
$2\arcsin\sqrt{\min(2r/L,1)}/\pi$, and
$\arcsin u\le2u$ for $0\le u\le1$. Layer-cake integration then gives
\[
 \int U^\rho\,d\omega_i-U^\rho(t_i)
 \le\sum_{j=1}^{16}\frac{8c_j\sqrt{2\varepsilon}}{\pi\sqrt{b_j-a_j}}
 \le60\sqrt\varepsilon.
\]
The last step uses $b_j-a_j>1/225$, $\sum c_j=\lambda<1$,
$\pi>3$, and $\sqrt2<3/2$.

The hypotheses of Lemma~\ref{lem:zeroenergy} need checking. The positive
measure $\eta=\sigma+\rho$ has compact support. The positive part of the
logarithmic kernel is bounded on that support. Circle potentials are
bounded below by $\log\varepsilon$. For an arcsine measure on $[a,b]$,
its potential at a complex point $z$ is at least its potential at
$\operatorname{Re}z$, because $|z-u|\ge|\operatorname{Re}z-u|$.
The real arcsine potential is bounded below by $\log((b-a)/4)$.
Consequently the negative logarithmic integral of $\eta$ is uniformly
bounded at points of its support, and Tonelli gives
$\iint|\log|z-w||\,d\eta(z)d\eta(w)<\infty$.
Since $|\sigma-\rho|\le\eta$ and both measures have mass $\lambda$,
Lemma~\ref{lem:zeroenergy} applies to $\sigma-\rho$.

Applying the circle-average identity successively yields
\[
 \iint\log|z-w|\,d\omega_i(z)d\omega_j(w)
 \ge\log\max(|t_i-t_j|,\varepsilon)\ge\log|t_i-t_j|.
\]
Therefore
$K^2I(\sigma)\ge2\sum_{i<j}\log|t_i-t_j|+h\log\varepsilon$.
Expanding $I(\sigma-\rho)\le0$, bounding the mixed terms by the
arcsine error, and applying \eqref{eq:potentialbound} gives exactly
\eqref{eq:configuration}.
\end{proof}

A uniform configuration bound alone cannot be integrated on an unbounded
domain. The next step explicitly leaves an integrable tail.

\begin{lemma}[Retaining confinement]
\label{lem:confinement}
For $K\ge2$ and $t\ge0$,
\[
 2U^\rho(t)-V(t)+\sqrt t/K\le M_0+\sqrt2/K.
\]
Consequently
\begin{equation}\label{eq:confinedconfiguration}
\begin{split}
 2\sum_{i<j}\log|t_i-t_j|-K\sum_iV(t_i)+\sum_i\sqrt{t_i}
 \le{}&(\lambda M_0-I(\rho))K^2\\
 &+(120+\sqrt2)h+2h\log K.
\end{split}
\end{equation}
\gid{Zeta5.Analysis.confined_configuration_bound}
\end{lemma}
\begin{proof}
For $t\le2$ use \eqref{eq:potentialbound}. For $t\ge2$, positivity and
support of $\rho$ imply $U^\rho(t)\le\lambda\log t$. Also
$\log(t+u^2)\le\log t+u^2/t$. Thus
\begin{equation}\label{eq:fieldtail}
 2U^\rho(t)-V(t)
 \le\frac{13}{10}\log t+\frac{2\alpha^3}{t}-2\pi\sqrt t.
\end{equation}
After adding $\sqrt t/K$, this right side is decreasing for $t\ge2$
and is less than $M_0$ at $2$. Direct differentiation and the elementary
bounds $\log2<7/10$, $\pi>3$, $\sqrt2>7/5$ suffice. Apply the proof of
Lemma~\ref{lem:regularization} to the field $V(t)-\sqrt t/K$, using
this bound and $\varepsilon=K^{-2}$. Since $h=\lambda K$, the
regularization error becomes $120h$ and the self-energy term becomes
$2h\log K$, which gives \eqref{eq:confinedconfiguration}.
\end{proof}

\subsection{The determinant integral and all scaling factors}
The finite determinant expansion and Fubini give
\begin{equation}\label{eq:Andreief}
 \Delta_K(\xi)=\frac1{h!}\int_{(0,\infty)^h}
 \prod_{i<j}(y_i^2-y_j^2)^2
 \prod_{i=1}^h\frac{D_N(y_i^2)^6}{D_K(y_i^2)}w(y_i)\,dy_i.
\end{equation}
To prove this identity without importing a named determinant integration
theorem, expand both Vandermonde determinants. The $h!$ choices of one
permutation give identical moment-determinant expansions after relabeling
integration variables. Every separated monomial integral is absolutely
convergent by the exponential decay of $w$, which justifies Fubini and
the finite rearrangements. Thus \eqref{eq:Andreief} is the required
instance of Andr\'eief's identity, with all its integrability hypotheses.
\gid{Zeta5.Analysis.gram_determinant_integral}

The explicit weight bound
\begin{equation}\label{eq:weightupper}
 w(y)\le8192(1+y)^5e^{-2\pi y}\qquad(y>0)
\end{equation}
follows from the rational generating function for $\sum\ell^4q^\ell$
and $1-e^{-u}\ge u/(1+u)$. Indeed these give
$w(y)\le\pi^{-1}(1+2\pi y)^5e^{-2\pi y}$, and $\pi<4$ supplies
$32\pi^4<8192$.

For $1\le m\le K$ and $t\ge0$, monotonicity gives
\begin{equation}\label{eq:riemannlog}
 0\le\sum_{j=1}^m\log\bigl(t+(j/K)^2\bigr)
    -K\int_0^{m/K}\log(t+u^2)\,du\le2\log K+2.
\end{equation}
For $t>0$, differentiating the difference in $t$ gives a nonpositive
quantity because $u\mapsto(t+u^2)^{-1}$ is decreasing. At $t=0$, the
upper bound follows from
$\log(m!)\le m\log m-m+\log m+1$. The continuous limit at zero
justifies including that endpoint.

In \eqref{eq:Andreief} set $y_i=K\sqrt{t_i}$. The separate powers of
$K$, before the Riemann-sum error, are
\[
\begin{array}{rl}
 K^{2h(h-1)}&\text{from the Vandermonde},\\
 K^{(12N-2K)h}&\text{from }D_N^6/D_K,\\
 K^{6h}&\text{from the weight and Jacobian}.
\end{array}
\]
The numerator error in \eqref{eq:riemannlog} adds at most
$K^{12h}e^{12h}$; the denominator uses the lower Riemann-sum bound.
Consequently
\begin{equation}\label{eq:scaledgram}
\begin{split}
 \Delta_K(\xi)\le{}&
 \frac{K^{2h(h+6N-K)+16h}(4096e^{12})^h}{h!}\\
 &\mathrel{}\times\int_{(0,\infty)^h}\prod_{i<j}(t_i-t_j)^2
       \prod_i e^{-KV(t_i)}t_i^{-1/2}(1+\sqrt{t_i})^5\,dt_i.
\end{split}
\end{equation}
There is no appeal to a uniform-bound-times-infinite-volume argument:
use \eqref{eq:confinedconfiguration} and retain
$\prod_i e^{-\sqrt{t_i}}$ in the integrand. Since
\[
 \int_0^\infty t^{-1/2}(1+\sqrt t)^5e^{-\sqrt t}\,dt
 =2\sum_{j=0}^5\binom5j j!=652,
\]
we obtain
\begin{equation}\label{eq:logdelta}
 \log\Delta_K(\xi)
 \le2h(h+6N-K)\log K+(\lambda M_0-I(\rho))K^2
                         +18h\log K+160h.
\end{equation}
The linear constant bounds
$\log4096+12+\log652+120+\sqrt2$ from above, and $h!\ge1$.

Elementary factorial bounds and summation give
\begin{align*}
 m\log m-m+1&\le\log(m!)\le m\log m-m+\log m+2,\\
 -2\sum_{i=1}^{h-1}\log((2i)!)
 &\le-2h^2\log(2h)+3h^2+4h\log(2h).
\end{align*}
Using \eqref{eq:scalar} yields
\begin{equation}\label{eq:logscalar}
 \log S_K\le(2\lambda-12\alpha\lambda-2\lambda^2)K^2\log K
                +C_*K^2+6h\log K+6h.
\end{equation}
The leading coefficient of $K^2\log K$ in \eqref{eq:logdelta} is
$2\lambda(\lambda+6\alpha-1)$, exactly the negative of the coefficient
in \eqref{eq:logscalar}.

\begin{proposition}[Real bound; source Proposition 6.3]
\label{prop:realbound}
For every $K\in40\Z_{>0}$,
\begin{equation}\label{eq:realbound}
 F_K(\xi)>0,\qquad
 \log F_K(\xi)\le UK^2+24K\log K+200K.
\end{equation}
\gid{Zeta5.real_determinant_bound}
\end{proposition}
\begin{proof}
Positivity is Proposition~\ref{prop:moments} and $S_K>0$. Add
\eqref{eq:logdelta} and \eqref{eq:logscalar}, cancel the leading
logarithmic terms, and use \eqref{eq:energyconstant} and $h\le K$.
Enlarging the remaining linear coefficient to $200K$ gives the stated
bound.
\end{proof}

\section{Completion of the polynomial estimates}
\label{sec:finish}
\begin{proof}[Proof of Theorem~\ref{thm:polynomialtarget}]
Proposition~\ref{prop:integrality} proves integrality for every
admissible pair. Proposition~\ref{prop:degree} proves degree $h$.
Proposition~\ref{prop:moments} and positivity of both rational
normalizations prove positive evaluation.

For fixed $M\in40\Z_{>0}$, \eqref{eq:AMbound} and
Proposition~\ref{prop:realbound} imply
\begin{equation}\label{eq:finalasymptotic}
 \limsup_{\substack{K\to\infty\\40\mid K}}
        K^{-2}\log Q_{K,M}(\xi)\le A_M+U.
\end{equation}
Appendix~\ref{app:arithmetic} proves, with exact positive margins,
\[
 -1600(A_{200}+U)>\frac{139}{5},\qquad
 -1600(A_{100000}+U)>\frac{7907}{100}.
\]
Choose a positive slack strictly smaller than each margin after division
by $1600$. The definition of the limsup bound then gives a threshold
for which the corresponding strict inequality holds for every allowed
$K$. Increase that threshold to satisfy $K\ge200M^2$, and put $K=40n$.
Since $K^2=1600n^2$, these are exactly
\eqref{eq:decay200}--\eqref{eq:decaylarge}. This establishes the theorem
with existential, not numerically asserted, decay thresholds.
\end{proof}
Theorem~\ref{thm:main} now follows from the integer-polynomial criterion.
No height estimate and no information about rational roots near $\xi$
are needed for this main endpoint.

\appendix
\section{The comparison measure and a reproducible finite certificate}
\label{app:measure}
\subsection{Rational data and the arcsine formula}
Let $\omega_{[a,b]}$ be the probability measure on $(a,b)$ with density
$1/(\pi\sqrt{(t-a)(b-t)})$, and put
\[
 \rho=\sum_{j=1}^{16}c_j\omega_{[a_j,b_j]}.
\]
Table~\ref{tab:arcsine} specifies the data by integers with common
denominator $10^{12}$. Exact arithmetic gives
\[
 \sum_jc_j=37/40,\qquad
 0<a_{16}<\cdots<a_1<b_1<\cdots<b_{16}<2,\qquad
 b_j-a_j>1/225.
\]
\begin{table}[htbp]
\centering\small
\caption{The rational arcsine measure (source Table 1).}\label{tab:arcsine}
\begin{tabular}{rrrr}\toprule
$j$ & $10^{12}a_j$ & $10^{12}b_j$ & $10^{12}c_j$\\\midrule
1 & 3906748086 & 8992695531 & 10515596180\\
2 & 2312248264 & 15340997855 & 29471737793\\
3 & 1402286665 & 25730180724 & 42934365099\\
4 & 881725356 & 41909578246 & 58204231966\\
5 & 578197906 & 65851089563 & 69037621310\\
6 & 396324613 & 99481037884 & 78873099189\\
7 & 283911191 & 144325727458 & 84856120711\\
8 & 212206188 & 201105762729 & 88396082127\\
9 & 165097686 & 269345996903 & 88303382125\\
10 & 133347132 & 347089554156 & 85472321255\\
11 & 111522114 & 430806704415 & 78899184238\\
12 & 96349355 & 515561896511 & 70353471918\\
13 & 85815639 & 595448778546 & 58838976615\\
14 & 78667711 & 664241383483 & 44421321106\\
15 & 74129565 & 716160577112 & 30462865791\\
16 & 71741310 & 746637295669 & 5959622577\\
\bottomrule\end{tabular}
\end{table}
The unit arcsine potential, with the present sign convention, is
\begin{equation}\label{eq:arcsinepotential}
 U^{\omega_{[a,b]}}(t)=
 \begin{cases}
 \log((b-a)/4),&a\le t\le b,\\
 \displaystyle\log\left(
 \frac{|t-(a+b)/2|+\sqrt{(t-a)(t-b)}}2\right),&t\notin[a,b].
 \end{cases}
\end{equation}
In particular the square root and absolute value are inside the
logarithm as displayed. To derive it, put $m=(a+b)/2$, $r=(b-a)/2$,
and substitute $u=m+r\cos\theta$. For $t>b$, differentiation of the
integral gives $((t-a)(t-b))^{-1/2}$; integrate and use
$U(t)-\log t\to0$ at infinity. Reflection treats $t<a$. For
$t=m+r\cos\phi$ in the interval use
$\cos\phi-\cos\theta=-2\sin((\phi+\theta)/2)\sin((\phi-\theta)/2)$.
The integral of the logarithms gives $\log(r/2)$, and continuity includes
the endpoints.

With $S_j=\sum_{i\le j}c_i$ and $S_0=0$, nesting gives
\begin{equation}\label{eq:energyexplicit}
 I(\rho)=\sum_{j=1}^{16}(S_j^2-S_{j-1}^2)
                            \log\frac{b_j-a_j}{4}.
\end{equation}
Indeed, the potential of component $j$ is constant on every earlier
support. Group its self-energy and both orders of its mixed energies;
the coefficient is $c_j^2+2c_jS_{j-1}=S_j^2-S_{j-1}^2$.

\subsection{Elementary interval arithmetic and its soundness contract}
For $1\le x\le2$, $z=(x-1)/(x+1)$, and $m\ge1$,
\begin{equation}\label{eq:logcertificate}
 0\le\log x-2\sum_{k=0}^{m-1}\frac{z^{2k+1}}{2k+1}
 \le\frac{2z^{2m+1}}{(2m+1)(1-z^2)}.
\end{equation}
For general positive $x$, write $x=2^e r$ with $1\le r<2$ and use
$\log x=e\log2+\log r$. For $0\le z\le1/2$,
\begin{equation}\label{eq:atancertificate}
 \left|\arctan z-\sum_{k=0}^{m-1}(-1)^k\frac{z^{2k+1}}{2k+1}\right|
 \le\frac{z^{2m+1}}{2m+1}.
\end{equation}
The reductions are
\[
 \arctan x=\frac\pi2-\arctan(1/x)\quad(x>1),\qquad
 \arctan x=2\arctan\frac{x}{1+\sqrt{1+x^2}}\quad(0\le x\le1),
